% Gesamttest «Trigonometrische Berechnungen» — Quelle des PDFs (python3 scripts/build-lp-pdf.py trigonometrische-berechnungen).
% Musterlösung und Raster: bewertungspaket.tex. Alle Werte mit python3 nachgerechnet (07.10.2026).
% Kein \lt/\gt — pdfLaTeX kennt sie nicht, < und > tun es.
\documentclass[11pt]{article}
\usepackage{../lp-druck}
\renewcommand{\lpname}{Trigonometrische Berechnungen}
\renewcommand{\lpteilgebiet}{5.3}
\renewcommand{\lpbereich}{Grundlagenfach}
\renewcommand{\lpfuss}{Gesamttest · 25 Punkte}
\begin{document}
\lpkopf{Gesamttest}{%
  \vspace{4pt}\noindent\begin{tabularx}{\linewidth}{@{}X X@{}}
  Code (kein Name): \hrulefill & Punkte: \hrulefill\ / 25
  \end{tabularx}}

\begin{lpinfo}{Anleitung}
\textbf{Zeit:} rund 30 Minuten. \textbf{Hilfsmittel:} Taschenrechner erlaubt, im Gradmodus (RLP 5.3 trägt keinen Vermerk «auch ohne Hilfsmittel»).
Mach zu jeder Sachaufgabe eine \textbf{Skizze} und schreib den \textbf{Rechenweg} auf — bewertet wird der Weg. Längen und Winkel
auf zwei Dezimalen runden, Zwischenresultate ungerundet weiterverwenden. Figuren sind Skizzen, nicht massstäblich.
Danach: Lösung scannen oder fotografieren (ohne Namen und Standort) und mit dem \emph{Bewertungspaket} bewerten lassen.
\end{lpinfo}

\lpteil{Teil A · Rechtwinklige Dreiecke}{Kapitel 1 bis 3 · 11 Punkte}

\begin{lpaufgabe}{G1}{Seiten und Winkel}{4 P}
\begin{minipage}[t]{0.55\linewidth}\vspace{0pt}
Im Dreieck $PQR$ liegt der rechte Winkel bei $Q$. Es ist $\angle P = 36^\circ$ und $\overline{PR} = 11\,\text{cm}$.
\begin{enumerate}[label=(\alph*),leftmargin=*,itemsep=2pt]
\item Welche Seite ist die Gegenkathete, welche die Ankathete des Winkels bei $P$?
\item Berechne $\overline{PQ}$ und $\overline{QR}$.
\item Wie gross ist der Winkel bei $R$?
\end{enumerate}
\end{minipage}\hfill
\begin{minipage}[t]{0.40\linewidth}\vspace{0pt}\centering
\begin{tikzpicture}[scale=0.42]
  \coordinate (P) at (0,0); \coordinate (Q) at (8.9,0); \coordinate (R) at (8.9,6.47);
  \draw[lpblau,thick,fill=lpblau!8] (P)--(Q)--(R)--cycle;
  \draw (8.9,0) ++(-0.5,0) -- ++(0,0.5) -- ++(0.5,0);
  \draw[lporange,thick] (1.4,0) arc (0:36:1.4); \node[lporange] at (2.1,0.55) {\small $36^\circ$};
  \node[below left] at (P) {$P$}; \node[below right] at (Q) {$Q$}; \node[above right] at (R) {$R$};
  \node[above left] at (4.45,3.3) {\small $11\,\text{cm}$};
\end{tikzpicture}
\end{minipage}
\lpplatz{4}
\end{lpaufgabe}

\begin{lpaufgabe}{G2}{Satteldach}{3 P}
Ein symmetrisches Satteldach ist $10\,\text{m}$ breit (Traufe bis Traufe), der First liegt $3.2\,\text{m}$ über der Traufe.
\begin{enumerate}[label=(\alph*),leftmargin=*,itemsep=2pt]
\item Wie gross ist der Neigungswinkel des Dachs?
\item Wie lang ist ein Sparren von der Traufe bis zum First?
\end{enumerate}
\lpplatz{4}
\end{lpaufgabe}

\begin{lpaufgabe}{G3}{Zwei Häuser vom Turm aus}{4 P}
Von der Plattform eines Aussichtsturms, $32\,\text{m}$ über dem Boden, siehst du zwei Häuser, die in derselben Richtung
hintereinander liegen: das nähere unter dem Tiefenwinkel $24^\circ$, das fernere unter $15^\circ$. Der Boden ist eben.
Wie weit liegen die beiden Häuser auseinander? Zeichne eine Skizze mit den beiden Tiefenwinkeln.
\lpplatz{6}
\end{lpaufgabe}

\lpteil{Teil B · Allgemeine Dreiecke}{Kapitel 4 und 5 · 14 Punkte}

\begin{lpaufgabe}{G4}{Waldbrand}{4 P}
Zwei Feuerwachen $A$ und $B$ liegen $6.5\,\text{km}$ auseinander. Ein Feuer $F$ wird von $A$ unter $\angle FAB = 48^\circ$ und
von $B$ unter $\angle FBA = 64^\circ$ gesehen.
\begin{enumerate}[label=(\alph*),leftmargin=*,itemsep=2pt]
\item Wie weit ist das Feuer von $A$ entfernt?
\item Wie weit ist das Feuer von der Verbindungslinie $AB$ entfernt?
\end{enumerate}
\lpplatz{5}
\end{lpaufgabe}

\begin{lpaufgabe}{G5}{Zwei mögliche Dreiecke}{3 P}
Gegeben sind $\alpha = 42^\circ$, $c = 10\,\text{cm}$ und die Gegenseite $a = 8\,\text{cm}$ von $\alpha$.
\begin{enumerate}[label=(\alph*),leftmargin=*,itemsep=2pt]
\item Begründe mit der Höhe von $B$ auf den Schenkel von $\alpha$, wie viele Dreiecke es gibt.
\item Berechne für das Dreieck mit dem \emph{stumpfen} Winkel $\gamma$ diesen Winkel und den Winkel $\beta$.
\end{enumerate}
\lpplatz{4}
\end{lpaufgabe}

\begin{lpaufgabe}{G6}{Grundstück}{4 P}
\begin{minipage}[t]{0.55\linewidth}\vspace{0pt}
Ein dreieckiges Grundstück hat zwei Seiten von $60\,\text{m}$ und $85\,\text{m}$, die den Winkel $125^\circ$ einschliessen.
\begin{enumerate}[label=(\alph*),leftmargin=*,itemsep=2pt]
\item Wie lang ist die dritte Seite?
\item Wie gross ist das Grundstück?
\end{enumerate}
\end{minipage}\hfill
\begin{minipage}[t]{0.40\linewidth}\vspace{0pt}\centering
\begin{tikzpicture}[scale=0.045]
  \coordinate (A) at (0,0); \coordinate (B) at (85,0); \coordinate (C) at (-34.41,49.15);
  \draw[lpblau,thick,fill=lpblau!8] (A)--(B)--(C)--cycle;
  \draw[lporange,thick] (8,0) arc (0:125:8); \node[lporange] at (5,11) {\small $125^\circ$};
  \node[below] at (42.5,0) {\small $85\,\text{m}$}; \node[left] at (-17,25) {\small $60\,\text{m}$};
\end{tikzpicture}
\end{minipage}
\lpplatz{5}
\end{lpaufgabe}

\begin{lpaufgabe}{G7}{Eine fremde Lösung}{3 P}
Lia soll im Dreieck mit $a = 5$, $b = 6$, $c = 10$ den Winkel $\gamma$ berechnen. Sie schreibt:
\[ \cos\gamma = \frac{25 + 36 - 100}{60} = -0.65 \]
«Ein Cosinus kann nicht negativ sein, also habe ich mich verrechnet.» Hat sie recht? Bestimme $\gamma$ und begründe mit einem Vergleich von $c^2$ und $a^2 + b^2$, was für ein Winkel $\gamma$ sein muss.
\lpplatz{3}
\end{lpaufgabe}
\end{document}
